\subsubsection{Direct measurement responses and fail-closed cache}
\label{sec:reduced-od-response-atoms}

Phase~6 constructs the observation operator directly from the complete journey
events in Section~\ref{sec:reduced-od-journey-choices}.  The implementation is
split between
\path{public_transportation.preprocessing.reduced_od.response_atoms},
\path{equivalence}, and \path{persistence}.  It never materializes or applies
the detailed OD-to-link assignment matrix.  Consequently this preprocessing
step does not import the assignment, sharding, microsharding, or
block-coordinate implementations.

Each supplied atomic boarding or alighting record is resolved against the
compact timetable arrays using its scenario stop, event time, trip or line, and
measurement type.  It must identify exactly one scheduled trip event.  Unknown,
unmatched, ambiguous, negative, or nonfinite records are rejected, and load
measurements remain unsupported by this first reduced backend.  A missing
sensor creates no measurement row, whereas a supplied numerical zero remains
an observed row with target zero.  Two different collection methods may map to
the same scheduled event, but they remain two distinct measurement rows and are
never silently summed.

Let $c$ index one journey OD/first-boarding-period cell, let $j\in J_c$ index
its fixed alternatives, and let $s_{j\mid c}$ be the Phase-5 initial share.  If
$n_{mj}$ is the number of boarding or alighting events of alternative $j$ that
match atomic measurement $m$, the direct response coefficient is
\[
B_{mc}=\sum_{j\in J_c}s_{j\mid c}n_{mj}.
\]
Only positive coefficients are stored, in canonical measurement-major COO
order.  For free journey demand $x$ and fixed positive demand $x^{\mathrm{fix}}$,
prediction is
\[
\widehat{\bm y}
=\bm b^{\mathrm{fix}}+B\bm x,
\qquad
b^{\mathrm{fix}}_m
=\sum_{c\in C_{\mathrm{fix}}}B_{mc}x^{\mathrm{fix}}_c.
\]
A frozen cell has no column in $B$.  A frozen positive value contributes only
to the fixed offset; a frozen zero contributes neither a free coordinate nor an
offset.  In particular, an unreachable structural-zero key is accepted without
a journey choice when its fixed value is zero; a positive fixed value without a
feasible journey is rejected.  Thus estimator dimension depends on free cells,
not on the size of the full OD table.

Free cells are placed in the same exact response-equivalence class only when
their ordered measurement indices and the bytes of every float64 coefficient
are identical.  Empty columns are therefore equivalent to one another, but an
approximately equal column is not compressed.  The artifact retains the class
of every free cell, its representative, and a CSR-style member list.  Phase~7
uses this partition to construct bounded operators without changing
predictions.

The compressed cache is written atomically as an NPZ archive without pickle
payloads.  Its metadata includes configuration, timetable, journey-choice, and
measurement fingerprints, canonical free/fixed keys, resolved atomic records,
and a fingerprint of every numerical array.  Loading reconstructs all immutable
contracts, recomputes the complete artifact fingerprint, and may require exact
expected upstream fingerprints.  Missing members, unsupported schema versions,
corrupt archives, modified contents, or stale identities fail closed.

The Phase-6 benchmark used a boarding and alighting record for every synthetic
stop-time event.  For 25, 50, 100, and 200 stops, respectively, 216, 432, 880,
and 1,776 measurement rows were processed in approximately 0.0037, 0.0071,
0.0137, and 0.0282 seconds.  The responses contained 58, 110, 110, and 110
nonzeros, while compressed cache files occupied approximately 5.3, 7.5, 11.4,
and 19.3 KB.  Every full-sensor response column in this fixture was distinct,
so its equivalence ratio was one.  A separate partial-sensor test verifies
compression of genuinely identical columns.  These measurements isolate
preprocessing and cache construction; they are not full-network thresholds.

